% Part: lambda-calculus
% Chapter: representability
% Section: currying

\documentclass[../../../include/open-logic-section]{subfiles}

\begin{document}

\olfileid[hi]{lam}{rep}{cur} 
\olsection{करीकरण}

कोई $\lambd$-अमूर्त $\lambd[x][M]$ एक तर्क वाले फलन को निरूपित करता है;
जब हम अनेक तर्क स्वीकार करने वाला फलन परिभाषित करना चाहते हैं, तब यह खासा
बड़ा प्रतिबंध है। ऐसा करने का एक उपाय $\lambd$-कलन को विस्तारित करके युग्म,
त्रिक आदि बनाने की अनुमति देना होगा; तब, मसलन, तीन-स्थानी फलन
$\lambd[x][M]$ अपने तर्क के त्रिक होने की अपेक्षा करेगा। किंतु इसे
\emph{करीकरण (Currying)} द्वारा करना अधिक सुविधाजनक है।

एक उदाहरण पर विचार करें। क्षण भर के लिए मान लें कि $\lambd$-कलन में हमारे
पास $+$ संक्रिया है। योग-फलन $2$-स्थानी है, अर्थात् वह दो तर्क लेता है।
किंतु कोई $\lambd$-अमूर्त हमें केवल एक तर्क वाले फलन देता है: वाक्यविन्यास
$\lambd[(x,y)][(x+y)]$ जैसे व्यंजकों की अनुमति नहीं देता। फिर भी, हम
$\lambd[y][(x+y)]$ द्वारा दिए एक-स्थानी फलन~$f_x(y)$ पर विचार कर सकते हैं,
जो अपने एकमात्र तर्क~$y$ में $x$ जोड़ता है। वास्तव में, यह एक अकेला फलन
नहीं बल्कि अलग-अलग ``$x$ जोड़ो'' फलनों का परिवार है---प्रत्येक संख्या~$x$
के लिए एक। अब हम दूसरा एक-स्थानी फलन~$g$, ~$\lambd[x][f_x]$ के रूप में
परिभाषित कर सकते हैं। तर्क $x$ पर लागू करने पर $g(x)$ फलन~$f_x$ लौटाता
है---अर्थात् उसके मान स्वयं अन्य फलन हैं। अब यदि हम $g$ को $x$ पर और फिर
परिणाम को~$y$ पर लागू करें, तो मिलता है: $(g(x))y = f_x(y) = x+y$। इस
प्रकार एक-स्थानी फलन~$g$, दो-स्थानी योग-फलन का ही काम कर सकता है।
``करीकरण'' इसी युक्ति का नाम है, जो दो-स्थानी फलनों को ऐसे एक-स्थानी फलनों
में बदलती है जिनके मान स्वयं एक-स्थानी फलन होते हैं।

अब $\lambd$-कलन के वाक्यविन्यास में ठीक से लिखा उदाहरण देखें। फलन
$f(x,y) = x$ को हम कैसे निरूपित करें? यदि हम ऐसा फलन परिभाषित करना चाहें
जो दो तर्क स्वीकार करके पहला लौटाए, तो $\lambd[x][\lambd[y][x]]$ लिख सकते
हैं। अक्षरशः यह ऐसा फलन है जो तर्क~$x$ स्वीकार करके फलन~$\lambd[y][x]$
लौटाता है। फलन $\lambd[y][x]$ दूसरा तर्क~$y$ स्वीकार तो करता है, पर उसे
छोड़ देता है और सदा~$x$ लौटाता है। देखें कि
$\lambd[x][\lambd[y][x]]$ को दो तर्कों पर लागू करने से क्या होता है:
\begin{align*}
  (\lambd[x][\lambd[y][x]])MN 
  \bredone & (\lambd[y][M])N \\
  \bredone & M
\end{align*}

सामान्यतः, किसी पद~$N$ द्वारा परिभाषित प्राचलों $x_1$, \dots,~$x_n$ वाला
फलन लिखने के लिए हम $\lambd[x_1][\lambd[x_2][\ldots\lambd[x_n][N]]]$
लिख सकते हैं। उस पर $n$ तर्क लागू करने से मिलता है:
\begin{multline*}
  (\lambd[x_1][\lambd[x_2][\ldots\lambd[x_n][N]]]) M_1 \dots M_n \bredone\\
  \begin{aligned}
  \bredone {} & (\Subst{(\lambd[x_2][\ldots\lambd[x_n][N]])}{M_1}{x_1}) M_2
  \dots M_n\\
   \eqs {} & (\lambd[x_2][\ldots\lambd[x_n][\Subst{N}{M_1}{x_1}]]) M_2
            \dots M_n \\
  \vdots & \\
  \bredone {} &\Subst{\Subst{P}{M_1}{x_1}\ldots}{M_n}{x_n}
  \end{aligned}
\end{multline*}
अंतिम पंक्ति का अक्षरशः अर्थ है फलन-परिभाषा के मुख्य भाग में $x_i$ के
स्थान पर $M_i$ प्रतिस्थापित करना; किसी फलन पर अनेक तर्क लागू करते समय हम
ठीक यही चाहते हैं।
\end{document}

